%%
%% main.tex — Example manuscript using ipparis-thesis class
%%
%% This file demonstrates the template capabilities.
%% Copy it and adapt it to your own thesis or HDR.
%%
%% Compile with: latexmk main.tex
%%           or: xelatex main && biber main && xelatex main && xelatex main
%%

\documentclass[phd,french,binding]{ipparis-thesis}
% Options: hdr | phd, french | english (default), cotutelle, binding

%% ── Metadata ─────────────────────────────────────────────────────

\author{Antoine Lavignotte}

\titlefr{Titre du manuscrit en fran\c{c}ais}
\titleen{Manuscript Title in English}

\doctoralschool{626}{Institut Polytechnique de Paris}{EDIPP}
\specialty{Computer Science, Data, AI}

\school{TSP}    % TSP, TP, X, ENSTA, ENSAE, PONTS
\lab{SAMOVAR}

\defensedate{\ifthesisfrench{1\textsuperscript{er} janvier 2026}{January 1, 2026}}
\defenseplace{Palaiseau}

\nnt{2026IPPAXXXX}

%% ── Jury ─────────────────────────────────────────────────────────

\ifthesisfrench{%
  \jurymember{Pr\'enom Nom}{Professeur, Universit\'e}{Rapporteur}
  \jurymember{Pr\'enom Nom}{Professeur, Universit\'e}{Rapporteur}
  \jurymember{Pr\'enom Nom}{Professeur, Universit\'e}{Rapporteur}
  \jurymember{Pr\'enom Nom}{Professeur, Universit\'e}{Examinateur}
  \jurymember{Pr\'enom Nom}{Professeur, Universit\'e}{Examinateur}
  \jurymember{Pr\'enom Nom}{Professeur, Universit\'e}{Pr\'esident}
}{%
  \jurymember{First Last}{Professor, University}{Reviewer}
  \jurymember{First Last}{Professor, University}{Reviewer}
  \jurymember{First Last}{Professor, University}{Reviewer}
  \jurymember{First Last}{Professor, University}{Examiner}
  \jurymember{First Last}{Professor, University}{Examiner}
  \jurymember{First Last}{Professor, University}{Chair}
}

%% ── Abstracts ────────────────────────────────────────────────────

\keywordsfr{mot-cl\'e 1, mot-cl\'e 2, mot-cl\'e 3}
\keywordsen{keyword 1, keyword 2, keyword 3}

\abstractfr{%
R\'esum\'e en fran\c{c}ais. Lorem ipsum dolor sit amet, consectetur adipiscing elit. Sed do eiusmod tempor incididunt ut labore et dolore magna aliqua.
}

\abstracten{%
Abstract in English. Lorem ipsum dolor sit amet, consectetur adipiscing elit. Sed do eiusmod tempor incididunt ut labore et dolore magna aliqua.
}

%% ── Bibliography ─────────────────────────────────────────────────

\addbibresource{bibliography.bib}

%% ── Glossary / Acronyms ─────────────────────────────────────────

\makeglossaries

\newacronym{ml}{ML}{Machine Learning}
\newacronym{rl}{RL}{Reinforcement Learning}

%% ═════════════════════════════════════════════════════════════════
\begin{document}
%% ═════════════════════════════════════════════════════════════════

%% ── Covers and front matter ──────────────────────────────────────

\frontmatter

% Official IP Paris front cover
\makefrontcover

% Dedication (optional)
\begin{dedication}
  \ifthesisfrench{\`A ceux qui liront ce document.}{To those who will read this document.}
\end{dedication}

% Acknowledgements
\begin{acknowledgements}
  \ifthesisfrench{Les remerciements vont ici.}{Acknowledgements go here.}
\end{acknowledgements}

% Table of contents
\tableofcontents

% Acronyms
\printglossary[type=\acronymtype,title=\ifthesisfrench{Acronymes}{Acronyms}]

%% ── Foreword ─────────────────────────────────────────────────────

\begin{foreword}
  \ifthesisfrench{%
    Ce document est un exemple du template \texttt{ipparis-thesis} pour la r\'edaction de manuscrits de th\`ese et d'HDR \`a l'Institut Polytechnique de Paris.

    Il illustre les fonctionnalit\'es disponibles~: couvertures officielles, environnements (d\'edicace, remerciements, avant-propos), chapitres, annexes et quatri\`eme de couverture.
  }{%
    This document is an example of the \texttt{ipparis-thesis} template for writing thesis and HDR manuscripts at Institut Polytechnique de Paris.

    It demonstrates the available features: official covers, environments (dedication, acknowledgements, foreword), chapters, appendices, and back cover.
  }
\end{foreword}

%% ── Main matter ──────────────────────────────────────────────────

\mainmatter

\chapter{\ifthesisfrench{Chapitre d'exemple}{Example chapter}}
\label{ch:example}

\ifthesisfrench{%
  Ceci est un chapitre d'exemple illustrant le template. Toutes les fonctionnalit\'es standard de \LaTeX{} sont utilisables~: \'equations, figures, tableaux, citations, r\'ef\'erences crois\'ees, etc.
}{%
  This is an example chapter demonstrating the template. You can use all standard \LaTeX{} features: equations, figures, tables, citations, cross-references, etc.
}

\section{\ifthesisfrench{\'Equations}{Equations}}

\ifthesisfrench{Un exemple d'\'equation~:}{An example equation:}
\begin{equation}
  E = mc^2
  \label{eq:einstein}
\end{equation}

\section{Citations}

\ifthesisfrench{%
  Citation d'une r\'ef\'erence~: voir~\cite{murphy2023survey} pour plus de d\'etails.
}{%
  Citing a reference: see~\cite{murphy2023survey} for details.
}

\section{\ifthesisfrench{Acronymes}{Acronyms}}

\ifthesisfrench{%
  Premi\`ere utilisation de \gls{ml} et \gls{rl}. Deuxi\`eme utilisation~: \gls{ml}.
}{%
  First use of \gls{ml} and \gls{rl}. Second use: \gls{ml}.
}

\chapter{\ifthesisfrench{Un autre chapitre}{Another chapter}}
\label{ch:another}

\ifthesisfrench{%
  Un autre chapitre pour illustrer la table des mati\`eres et les en-t\^etes de page.
}{%
  Another chapter to show the table of contents and page headers.
}

\section{\ifthesisfrench{Une section}{A section}}

\ifthesisfrench{Du contenu ici.}{Some content here.}

\subsection{\ifthesisfrench{Une sous-section}{A subsection}}

\ifthesisfrench{Encore du contenu.}{More content.}

%% ── Back matter ──────────────────────────────────────────────────

\backmatter

% Bibliography
\printbibliography[heading=bibintoc]

%% ── Appendices ───────────────────────────────────────────────────

\appendix

\chapter{\ifthesisfrench{Annexe d'exemple}{Example appendix}}
\label{app:example}

\ifthesisfrench{Le mat\'eriel compl\'ementaire va ici.}{Additional material goes here.}

%% ── Back cover ───────────────────────────────────────────────────

\makebackcover

\end{document}
